\documentclass[10pt,a4paper]{amsart}
\usepackage[margin=19mm]{geometry}
\usepackage{amsmath,amssymb,mathtools}
\usepackage[scr=rsfs,cal=euler]{mathalfa}
\usepackage{xfrac}
\usepackage[unicode,hidelinks,bookmarks=false]{hyperref}
\newcommand{\ie}{i.e.\ }
\newcommand{\cf}{cf.\ }
\newcommand{\resp}{respectively\ }
\newcommand{\Subsection}{\subsection}
\providecommand{\nobreakdash}{\nobreak\hbox{-}}
\setlength{\emergencystretch}{2em}
\multlinegap0pt
\usepackage{amssymb}
\usepackage{mathtools}
\usepackage{xcolor}
\usepackage[shortlabels]{enumitem}
\usepackage{bm}
\usepackage{tikz}
\usepackage{cancel}
\usepackage{siunitx}
\newtheorem{lemma}{Lemma}
\title{Some reducible and irreducible Brill--Noether loci.}
\author{Richard Haburcak, Montserrat Teixidor i Bigas}
\date{Source: arXiv:2503.16255v2 (CC BY 4.0)}
\begin{document}
\maketitle

\noindent\emph{Source and licence.} Some reducible and irreducible Brill--Noether loci.. Version arXiv:2503.16255v2 (CC BY 4.0), distributed by arXiv under the Creative Commons Attribution 4.0 International licence, http://creativecommons.org/licenses/by/4.0/, as recorded in the arXiv licence metadata for this exact version and shown on the abstract page https://arxiv.org/abs/2503.16255v2. Full licence notice: \texttt{licenses/arxiv-2503.16255-CC-BY-4.0.txt}.\par\smallskip

\noindent\textit{Editorial scope.} Counted unit: the article's lemma lemma (source0.tex lines 1309--1312). The article's complete original proof of it is delivered with it (source0.tex lines 1313--1320, one \texttt{proof} environment of 8 lines). No mathematics is written, repaired or re-derived: the statement and the proof are the article's own lines, in the article's own notation, and no retained line is substituted or re-flowed.  No further source-local context is needed: every cross-reference of every retained line resolves inside the two delivered spans.  Nothing was omitted from any retained span, so the package carries no omission record.  No retained line contains a citation, so the wrapper carries no bibliography and no bibliography entry.  No retained line contains a figure environment, an image-inclusion command or drawing code, so no image is delivered and the package declares no asset dependency.  Editorial formatting only, in addition: an AMS \texttt{amsart} preamble that restates the article's own declarations and macros used by the retained lines, namely the environments \texttt{lemma} on separate counters, together with the standard packages those lines need.  The retained environments print in this document as Lemma 1 in order of appearance, whereas the article prints the same items with the numbering of its own full document.  The complete native source of the article as distributed in the arXiv e-print for this exact version is bundled unchanged as \texttt{sources/175113/source0.tex}.
\begin{lemma}\label[lemma]{lemgon}
	Let $C$ be a chain of elliptic curves that is in the closure of the locus of $k$-gonal curves. 
	Then, in at least $-\rho$ of the elliptic components the points $P_i, Q_i$ satisfy $m_i(P_i-Q_i)=0$ and these $m_i$ necessarily satisfy  $m_i\le k$.
\end{lemma}

\begin{proof} If a chain of elliptic curves is the limit of $k$-gonal curves, then one can fill a $2\times (g-k+1)$ rectangle with the indices $1,2,\dots, g$.
	This requires a minimum of $2\times (g-k+1)-g$ repeats of indices. 
	Because one of the dimensions of the rectangle is 2, an index cannot appear more than twice.
	There are therefore at most $g-[2\times (g-k+1)-g]=2k-2$ non repeated indices in the filling of the rectangle.
	To obtain maximal torsion among the points $P_i, Q_i$, the two appearances of the index $i$ should be placed at the maximum possible distance.
	This can be achieved by placing the first $k-1$ indices in the first column to create a gap of $k$ then repeated indices one on each column and the last $k-1$ indices in the last column to close the gap.
	Note that if there are fewer non-repeated indices, the maximum $m_i$ will decrease rather than increase.
\end{proof}

\end{document}
